\documentclass[10pt, letterpaper]{article}

% =========================================================
% Page and Format Settings
% =========================================================
\usepackage{geometry}
\geometry{left=1.4cm, right=1.4cm, top=0.8cm, bottom=0.8cm}

\usepackage{enumitem}
\usepackage{titlesec}
\usepackage{array}
\usepackage{tabularx}
\usepackage[hidelinks]{hyperref}

\pagestyle{empty}
\setlength{\parindent}{0pt}
\renewcommand{\baselinestretch}{0.9}

% =========================================================
% Heading macros
% =========================================================
% Line 1: Organization / Project (bold) ...... Location
% Line 2: Role / Context (italic) ............ Dates
\newcommand{\entry}[4]{%
    \noindent
    \begin{tabularx}{\textwidth}{@{}>{\raggedright\arraybackslash}X r@{}}
        \textbf{#1} & #2 \\
        \textit{#3} & \textit{#4}
    \end{tabularx}\par\vspace{-1pt}%
}

% Institution line for grouped projects
\newcommand{\inst}[2]{%
    \noindent
    \begin{tabularx}{\textwidth}{@{}>{\raggedright\arraybackslash}X r@{}}
        \textbf{#1} & #2
    \end{tabularx}\par\vspace{-1pt}%
}
% Project under an institution: Title (bold italic) + note ......... Dates
\newcommand{\proj}[3]{%
    \noindent
    \begin{tabularx}{\textwidth}{@{}>{\raggedright\arraybackslash}X r@{}}
        \hspace{0.6em}\textbf{\textit{#1}}\ \textit{#2} & \textit{#3}
    \end{tabularx}\par\vspace{-1pt}%
}

\titleformat{\section}{\large\bfseries\scshape}{}{0em}{}[\titlerule]
\titlespacing{\section}{0pt}{5pt}{3pt}

\setlist[itemize]{noitemsep, topsep=1pt, parsep=0pt, partopsep=0pt,
    leftmargin=1.3em, after=\vspace{3pt}}

% =========================================================
\begin{document}

% ---------------- Header ----------------
% Public web version: phone number omitted.
\begin{center}
    {\LARGE \textbf{Haoxin Guo}} \\[3pt]
    Urbana, IL \ $|$ \
    \href{mailto:s611guo@gmail.com}{s611guo@gmail.com} \ $|$ \
    \href{https://www.linkedin.com/in/haoxin-guo-4a8431351}{linkedin.com/in/haoxin-guo-4a8431351} \ $|$ \
    \href{https://github.com/zndwb}{github.com/zndwb} \\[2pt]
    \textbf{Seeking Summer 2027 Internship \& Fall 2027 Co-op: Machine Learning Engineer}
\end{center}

% ---------------- Education ----------------
\section{Education}

\entry{University of Illinois Urbana-Champaign}{Urbana, IL}
      {Master of Computer Science}{Aug 2026 -- Expected May 2028}
\begin{itemize}
    \item Coursework: Advanced Topics in NLP (CS 546), Applied Machine Learning (CS 441), Web Programming (CS 409)
\end{itemize}

\entry{China Agricultural University}{Beijing, China}
      {B.E. in Computer Science and Technology; GPA: 3.67/4.00}{Sep 2022 -- Jun 2026}
\vspace{3pt}

% ---------------- Experience ----------------
\section{Internship Experience}

\entry{Illinois Institute of Technology}{Remote}
      {Research Intern (Advisor: Prof. Ren Wang)}{Jul 2025 -- Present}
\begin{itemize}
    \item \textbf{MoE unlearning (first author, under review at \textit{IEEE TPAMI}):} Showed that single-expert unlearning
          in \textbf{Mixture-of-Experts} LLMs is superficial (knowledge rebounds after brief fine-tuning)
          and proposed \textbf{PLEU}, a layer-wise expert erasure method. On DeepSeek-V2-Lite and
          Qwen1.5-MoE, it had \textbf{lower knowledge recovery} than single-expert baselines in every setting,
          preserving MMLU utility with \textbf{67--85\% less} runtime than mode-connectivity unlearning.
    \item \textbf{Agent unlearning:} Core contributor to a project on unlearning for \textbf{self-evolving LLM agents},
          removing knowledge stored in agent memory and experience, not only weights; designed the experimental setup
          and ran experiments.
\end{itemize}

\entry{DP Technology}{Beijing, China}
      {Machine Learning Engineer Intern}{Oct 2025 -- Jan 2026}
\begin{itemize}
    \item Pre-trained a \textbf{diffusion-based foundation model} from scratch on nanopore (POD5) current signals and
          re-engineered the PyTorch feature-extraction pipeline, reaching 93\% accuracy on core downstream benchmarks.
    \item Applied \textbf{graph representation learning} to RNA--RNA interaction prediction (97\% accuracy on in-house data).
    \item Containerized backend computation with \textbf{Docker} on the Bohrium platform and independently shipped a protein
          normal-mode-analysis web app with automated structure generation and visualization.
\end{itemize}

\entry{Datang Software Technology}{Beijing, China}
      {Software Engineer Intern}{Aug 2024 -- Nov 2024}
\begin{itemize}
    \item Developed Django backend services integrating an LLM API (Zhipu GLM) to automate text-generation workflows.
    \item Built data cleaning and ETL jobs with SQL and Spark, and tuned the platform database for query performance.
\end{itemize}

% ---------------- Projects ----------------
\section{Project Experience}

\inst{University of Illinois Urbana-Champaign}{Urbana, IL}
\proj{CulminationNLI: Event-Completion Reasoning in Multilingual LLMs}{(CS 546)}{Fall 2026 -- Present}
\begin{itemize}[leftmargin=2em]
    \item Co-developing an \textbf{English--Mandarin NLI benchmark} ($\sim$780 pairs) testing if LLMs tell event
          occurrence from completion, with labels validated by native speakers.
    \item Evaluating Qwen, Gemma, and Ministral under zero-shot, definition-aware, few-shot, and CoT prompting, with
          minimal-pair accuracy and over-/under-inference error analysis; paper planned for submission.
\end{itemize}

\inst{China Agricultural University}{Beijing, China}
\proj{Lightweight Tomato Disease Segmentation}{(first author)}{Jul 2023 -- Jan 2025}
\begin{itemize}[leftmargin=2em]
    \item Proposed a \textbf{MaxMin-Diffusion attention} mechanism with temporal segmentation, reaching 92\% accuracy and
          0.91 mIoU (+8 pts over CBAM, +19 over self-attention); deployed on-device on a smartphone.
\end{itemize}
\proj{Deep RL Framework for Automatic Hyperparameter Search}{(co-first author)}{Sep 2023 -- May 2025}
\begin{itemize}[leftmargin=2em]
    \item Led design and implementation of \textbf{DWIDRL}, framing hyperparameter search as an MDP with progressive
          search-space attenuation; tuning 5 trajectory-classification models across 4 harvester/tractor scenarios, it beat
          grid search, Bayesian optimization, PSO, and Q-learning by up to \textbf{+21\% F1} and \textbf{+15\% accuracy}.
\end{itemize}

% ---------------- Publications ----------------
\section{Publications}
\begin{itemize}[leftmargin=1.3em]
    \item \textbf{H. Guo}, Y. Shi, R. Wang. ``Beyond Superficial Unlearning: Progressive Layer-wise Expert Erasure in
          MoE Models.'' \textit{IEEE Transactions on Pattern Analysis and Machine Intelligence (TPAMI)} (IF 20.4), under review.
    \item \textbf{H. Guo}, et al. ``Efficient Tomato Disease Detection Using MaxMin-Diffusion Mechanism and Lightweight
          Techniques.'' \textit{Plants} (IF 4.7), 14(3):354, 2025. \textit{(First author)}
    \item W. Zhai$^*$, \textbf{H. Guo}$^*$, J. Pan, C. Wu. ``An Automatic Parameter Search Framework for Agricultural
          Machinery Trajectory Operation Mode Identification by Deeply Weakly Informed Deep Reinforcement Learning.''
          \textit{Computers and Electronics in Agriculture} (IF 10.3), 253:112172, 2026. \textit{($^*$Equal contribution)}
\end{itemize}

% ---------------- Skills ----------------
\section{Technical Skills}
\begin{itemize}[leftmargin=1.3em]
    \item \textbf{Languages:} Python, C/C++, SQL, MATLAB, R, Verilog, Bash
    \item \textbf{ML \& LLMs:} PyTorch, Hugging Face, LoRA/PEFT, SFT, RLHF/DPO, MoE, Diffusion Models, Graph Learning, RL
    \item \textbf{Systems \& Tools:} Docker, Linux, Git, Django, Spark, LangChain, RAG, LLM Agents
\end{itemize}

% ---------------- Awards ----------------
\section{Awards}
\begin{itemize}[leftmargin=1.3em]
    \item Grand Prize, Capital ``Challenge Cup''; Grand Prize, ``BeiDou Star'' National Finals; Honorable Mention, MCM
\end{itemize}

\end{document}
